\begin{annexitem}{annex:p05-rep-bounded-decisions}{thm:p05-rep-bounded-decisions}{Finite arithmetic and bounded decisions without induction}
\begin{proof}
We construct finite \(\mathsf Q\) derivations by inductions in the
metatheory. None is an application of an induction axiom in \(\mathsf Q\).
For standard \(a,b\), unroll \(a+0=a\) and
\(a+Sb=S(a+b)\) exactly \(b\) times; similarly unroll
\(a\times0=0\) and \(a\times Sb=(a\times b)+a\).
These give proofs evaluating each operation on numerals. Recursing over
the finite syntax of a closed term gives a proof that it equals its value
numeral. Successor injectivity reduces equality of distinct numerals to
\(S^r0=0\) for some \(r>0\), contradicted by \(Sx\ne0\).
Equality reflexivity handles equal numerals, and \(\bottom\) is refutable.

For the source order \(u\leq\bar n\),
\cref{lem:q-bounded-cases} gives the finite-case implication. Conversely,
each \(\bar j\leq\bar n\), \(j\leq n\), has the explicit witness
\(\overline{n-j}\) in \(z+\bar j=\bar n\), justified by the foregoing
evaluation. Equality substitution and disjunction elimination therefore
give both finite-quantifier equivalences for \(\leq\), for an arbitrary
formula matrix, not merely a bounded one.

The book's different orientation for \(<\) must be checked separately.
For each fixed \(n\), \(\mathsf Q\) proves
\[
 \forall x\forall z\bigl(x+Sz=\bar n\to
              \bigvee_{j<n}x=\bar j\bigr).
\]
For \(n=0\), \(x+Sz=S(x+z)\ne0\). For the passage from \(n\) to
\(n+1\), injectivity turns \(x+Sz=S\bar n\) into \(x+z=\bar n\).
Classical logic and the predecessor axiom give \(z=0\) or
\(z=Sw\). In the first case \(x=\bar n\). In the second, use the
already constructed proof for \(x+Sw=\bar n\). This is a finite
proof-construction recursion on the external \(n\). Conversely, for
each \(j<n\), the numeral witness \(\overline{n-j-1}\) proves
\(\bar j+S\overline{n-j-1}=\bar n\).
This proves the finite-quantifier equivalences for \(<\), including its
empty zero range, without asserting that its comparison is universally
\(\mathsf Q\)-equivalent to the other order abbreviation.

Finally perform structural recursion on a closed marked-bounded formula.
The atomic case was proved above; Boolean connectives use their classical
truth tables. In a bounded-quantifier case its closed bounding term is
first replaced by its numeral value. Replace the quantifier by the proved
finite conjunction or disjunction, substitute each numeral in the matrix,
and apply the already constructed proofs for those smaller formulas.
This also works when later bounds depend on earlier bound variables:
substitute the earlier fixed numeral before evaluating the later term.
Every branch is finite. It supplies the positive or negative derivation
according to standard truth, proving the assertion.
\end{proof}
\end{annexitem}

\begin{annexitem}{annex:p05-rep-trace}{thm:p05-rep-trace}{Compilation, beta packing, and exact trace checking}
\begin{proof}
We first give the promised program construction, including its finite
termination argument. A jump without changing registers is a zero-test
on a dedicated permanently zero register with identical targets.
A zeroing subroutine repeatedly tests/decrements its target register,
returning when it is zero. To copy a source \(A\) into a distinct target
\(B\) while preserving \(A\), first zero \(B\) and a fresh temporary
\(U\); repeatedly decrement \(A\), incrementing \(B,U\) each time;
then repeatedly decrement \(U\), incrementing \(A\) each time.
For initial standard value \(a\), the first loop after \(j\leq a\)
iterations has values \((A,B,U)=(a-j,j,j)\); the restoring loop after
\(j\leq a\) iterations has \((A,B,U)=(j,a,a-j)\).
Thus the loops halt after exactly \(a\) iterations apiece and return
\((a,a,0)\). Zeroing likewise terminates after the original target
value iterations. These are external finite invariants, not formal
\(\mathsf Q\) induction claims.

Construct subprograms by recursion on the finite description, with
specified input, output, scratch, and return-label registers. All input
registers are preserved, and scratch registers are fresh and cleared on
entry. For zero, clear the output; for a projection, copy the appropriate
input; for successor, copy its input then increment the output.
For composition \(g(h_1(\vec x),\ldots,h_m(\vec x))\), give each
\(h_j\) fresh input copies, scratch registers, and an output register.
Run these subprograms in order, copy their outputs into the dedicated
inputs of \(g\), and run \(g\). Relabel instruction addresses in each
copy and replace its halt by the next fixed return address. Inputs
\(\vec x\) are untouched. The already proved termination and correctness
of the finitely many subprograms prove termination and the composition
value. Nullary descriptions and empty argument lists use the same rule
with no input-copy operations.

For primitive recursion, preserve \((\vec x,n)\), copy \(n\) to a
counter \(C\), zero an index \(I\), and run \(g\) into an accumulator
\(A\). At the loop head test \(C\). If zero, copy \(A\) into the final
output and return. Otherwise decrement \(C\), give \(h\) fresh copies
of \((\vec x,I,A)\), run it, copy its output back to \(A\), increment
\(I\), and jump to the loop head. After \(j\leq n\) completed
iterations the invariant is
\(C=n-j, I=j, A=f_d(\vec x,j)\); the base is the defining equation
for \(g\), and the successor is the defining equation for \(h\).
Each invocation of \(h\) terminates by the description recursion.
There are exactly \(n\) iterations, followed by a zero exit, so the
output is \(f_d(\vec x,n)\). Fresh disjoint register blocks prevent
interference with \(C,I,A,\vec x\). There are only finitely many nested
descriptions, hence finitely many registers and instructions. These
recipes are effective finite text transformations. Put input registers
first and relabel the output as \(R_k\); unused registers start zero.
This proves the promised compilation for every description, not just a
termination assertion about an unspecified machine.

Next recall the exact sequence existence used to encode a standard run.
The numerical pairing satisfies its atomic graph \(\mathsf P\), and
\(\mathsf B(w,i,v)\) extracts the remainder of the uniquely decoded
first component by \(1+(i+1)m\), as proved in
\cref{lem:pairing-graph,lem:goedel-beta}. Given finitely many desired
entries, set the unspecified entries to zero up to the greatest index.
Choose \(N\) larger than that length and all entries, and \(m=N!\).
The moduli \(1+(i+1)m\) are pairwise coprime: a common prime divides
\((j-i)m\), not \(m\), hence \(j-i\); but \(0<j-i<N\) then makes it
a divisor of \(N!\), a contradiction.
For completeness, the finite Chinese-remainder step here is constructive:
for coprime \(a,b\), the Euclidean division algorithm terminates with gcd
\(1\); back-substitution gives integers \(u,v\) with \(ua+vb=1\).
Given residues \(r\pmod a,t\pmod b\), choose the nonnegative residue
\(h\) of \(u(t-r)\pmod b\). Then \(r+ah\) realizes both residues.
Its nonnegative residue modulo \(ab\) is a natural simultaneous
solution. Iterating this combination over the finite coprime list gives
the required \(t\), with all desired entries smaller than their moduli.
Pair \((t,m)\) to obtain \(w\). Euclidean division gives exactly the
desired beta entries, each at most \(w\). This is an argument about
standard finite sequences, not an internal assertion that all
nonstandard sequence requests have solutions in \(\mathsf Q\).

The two tagged extractors pair their indices, with the bound
\(\operatorname{pair}_{\N}(a,b)<(a+b+1)^2\), as proved in
\cref{lem:p04-tagged-history-extraction}. The arity of each family is
fixed; distinct positions give distinct paired indices. Therefore the
preceding finite packing realizes simultaneously the header, label and
register entries of any finite run in one \(s\). All assigned values,
including \(\ell\), the input and output, are at most \(s\).
The witnesses for \(\operatorname{Le}(y,s)\) can be chosen as \(s-y\),
which is also at most \(s\). No paired-index bound by \(s\) is assumed:
the extractor has its separate polynomial index bounds.

If the standard program run has output \(b\), pack this finite assignment.
Each transition satisfies its displayed increment or decrement clause;
unchanged registers satisfy their equality clauses; its initial and final
rows satisfy the other conjuncts. Thus \(C_d(\vec a,b,s)\) holds.
Conversely, suppose \(C_d(\vec a,b,s)\) holds in \(\mathcal N\).
Uniqueness of tagged beta extraction identifies all overlapping row
entries. At time zero they are exactly the input configuration. Induct
externally on \(i<\ell\). Its stored label is a valid instruction;
the finite step disjunction selects that instruction and gives its actual
deterministic successor. Thus the stored next row is the actual next
configuration. There is no step clause for a halt instruction, so no
earlier halt can be continued. The final stored halt and output give the
actual output \(b\), which by compilation is \(f_d(\vec a)\).
The zero-step case is also covered: it is accepted precisely if the fixed
program already halts with the specified initial output. There is no
hidden positive-length assumption.

All quantifiers in \(C_d\) are marked bounded, and the extractor
expansions are bounded; all other clauses are finite Boolean combinations
of equations. Fresh renaming avoids variable capture. To evaluate this
one fixed formula numerically, evaluate its terms primitively recursively,
use equality and Boolean tests, and evaluate each bounded existential by
a bounded sum followed by a positivity test, each bounded universal by a
bounded product. Substitute the primitively recursive value of the bound
into that recursion. Empty ranges give the correct Boolean values.
Structural recursion on the fixed formula proves its truth characteristic
primitive recursive. Finally existential elimination on
\(\exists w\leq s(w+y=s)\) immediately gives
\(\exists w(w+y=s)\), which is the source \(y\leq s\), already in
\(\mathsf Q\). This proves all assertions.
\end{proof}
\end{annexitem}

\begin{annexitem}{annex:p05-rep-strong-graph}{thm:p05-rep-strong-graph}{The least-certificate guard excludes nonstandard outputs}
\begin{proof}
Fix standard \(\vec a\) and \(m=f_d(\vec a)\). By standard trace
adequacy at least one certificate exists. Let \(r\) be the least
\emph{standard} number coding any successful trace for that input and
any output. Every successful standard trace has output \(m\), and its
output is at most its certificate. Thus \(m\leq r\),
\(C_d(\overline{\vec a},\bar m,\bar r)\) is true, and for every
standard \(t<r\),
\(\exists z\leq\bar t\,C_d(\overline{\vec a},z,\bar t)\) is false.
These are closed bounded facts. The bounded-decisions lemma proves their
truth values in \(\mathsf Q\). Finite aggregation over \(t\leq\bar r\)
then proves the minimum guard at \(\bar r\), using its equality
disjunct for \(t=\bar r\). Hence
\(\mathsf Q\derives\delta_d(\overline{\vec a},\bar m,\bar r)\),
and existential introduction proves \(\gamma_d\) at \(\bar m\).

We also need finite-certificate decisions with an unrestricted output.
For each standard \(t\leq r\), the logical consequence
\(C_d(\overline{\vec a},y,\bar t)\to y\leq\bar t\) reduces \(y\)
to its finitely many numeral cases. In each such case bounded decisions
refute the certificate when \(t<r\). For \(t=r\), they refute every
output other than \(\bar m\). Therefore \(\mathsf Q\) proves
\[
 \forall y\neg C_d(\overline{\vec a},y,\bar t)\ (t<r),\qquad
 \forall y(C_d(\overline{\vec a},y,\bar r)\to y=\bar m).
\]
These are actual finite derivations obtained by the bound on \(y\), not
an inference from all numeral instances to a universal statement.

Now argue in \(\mathsf Q\) from
\(\delta_d(\overline{\vec a},y,s)\), with arbitrary \(y,s\).
The fixed-standard comparison \cref{lem:q-dichotomy} gives
\(s\leq\bar r\lor\bar r\leq s\). In the first case finite cases
for \(s\) reduce it to \(\bar0,\ldots,\bar r\). The displayed
refutations leave only \(s=\bar r\), and then give \(y=\bar m\).
In the second case the premise is exactly
\(\exists u\,(u+\bar r=s)\), because the source defines
\(\bar r\leq s\) by that left-witness equation. Here is the complete
guard-instantiation derivation, including the relevant abbreviation
expansions. Put
\[
 J(t)\defeq\exists z\bigl((\exists v\,(v+z=t))\land
                                  C_d(\overline{\vec a},z,t)\bigr).
\]
Thus \(J(t)\) is precisely the expansion of
\(\exists z\leq t\,C_d(\overline{\vec a},z,t)\), not a version in
which the membership comparison has been replaced by
\(\operatorname{Le}(z,t)\). Fresh variables in each expansion avoid
capture. Conjunction elimination from the assumed \(\delta_d\) gives
\[
 \forall t\bigl((\exists u\,(u+t=s))\to(t=s\lor\neg J(t))\bigr).
\]
Universal elimination at the closed term \(\bar r\) gives
\[
 (\exists u\,(u+\bar r=s))\to(\bar r=s\lor\neg J(\bar r)).
\]
Implication elimination with the \emph{exact} second-case premise gives
\(\bar r=s\lor\neg J(\bar r)\). These three steps are first-order
logical derivations valid already over the axioms of \(\mathsf Q\);
they require no equation about arbitrary addition and do not require
\(\operatorname{Le}(\bar r,s)\).
For hierarchy bookkeeping, these expanded displays exhibit meaning and
logical inference, not a claimed bounded raw-language expansion.
\Cref{def:language-arithmetic} explicitly defines marked
\(\forall t\leq s\) by this implication, and declares such
\(\leq\)-marked quantifiers counted as bounded exactly like its
\(<\)-marked quantifiers. The following syntactic-hierarchy definition
generates bounded formulas from atomic formulas, Boolean connectives and
marked bounded quantifiers, and treats marked duality before eliminable
expansion. Thus the original displayed \(\delta_d\) is a bounded
matrix and \(\exists s\delta_d\) a literal existential-prefix formula
\emph{in that presentation}. The hidden unbounded comparison witnesses
in the full \(\Lang_A\) expansion have not been bounded by a
\(\mathsf Q\) equivalence. No such equivalence is used or needed here.

To eliminate its second disjunct, \(m\leq r\) is a relation between
fixed standard numbers. The already constructed finite addition proof is
\[
 \mathsf Q\derives
       \overline{r-m}+\bar m=\bar r.
\]
Existential introduction on its left witness gives
\(\exists v(v+\bar m=\bar r)\). Conjoin this with the previously
proved \(C_d(\overline{\vec a},\bar m,\bar r)\), then existentially
introduce \(z\) with witness \(\bar m\), obtaining \(J(\bar r)\).
Disjunction elimination now contradicts the \(\neg J(\bar r)\)
branch and leaves \(\bar r=s\). Equality symmetry gives \(s=\bar r\);
substitute it into \(C_d(\overline{\vec a},y,s)\) and use the proved
fixed-certificate universal implication to obtain \(y=\bar m\).
In particular the argument never infers a bounded witness from
\(\exists u(u+\bar r=s)\), never assumes \(0+s=s\), and never switches
to the different \(<\) comparison. The separate
\(\operatorname{Le}(y,s)\) conjunct in a certificate is used only in
the first, finite-certificate argument, where forgetting its witness bound
logically gives the source \(y\leq s\).
Both cases discharge to
\(\forall y\forall s(\delta_d(\overline{\vec a},y,s)\to y=\bar m)\).
Existential elimination gives the forward graph implication. The reverse
implication follows from the proved graph at \(\bar m\) by equality
substitution. Universal generalization on \(y\) proves the requested
biconditional.

In \(\mathcal N\), the minimum certificate and correct output satisfy
the graph, and any satisfying graph has a successful trace, whose output
is correct. Thus standard-model graph accuracy follows too. The only
comparison involving an arbitrary certificate used a fixed numeral
\(\bar r\). We used neither general order trichotomy nor induction on a
possibly nonstandard run nor universal existence of an internal minimum.
\end{proof}
\end{annexitem}

\begin{annexitem}{annex:p05-rep-numeralwise-relations}{thm:p05-rep-numeralwise-relations}{Characteristic graphs and bounded aggregates}
\begin{proof}
For any standard tuple, its characteristic value is either \(1\) or
\(0\). The strong-graph theorem gives
\(\forall y(\gamma_{d_R}(\overline{\vec a},y)\leftrightarrow y=\bar1)\)
in the first case and the corresponding biconditional with \(0\) in
the second. Specialize to \(y=\bar1\); reflexivity proves the positive
instance, and \(S0\ne0\) proves the negative instance. The literal
existential prefix and standard-model accuracy follow from the selected
definition and graph theorem. For a fixed finite bound, use the
\(\leq\)-finite-quantifier equivalence proved in the first annex item
with the actual matrix \(\rho_R\) or its negation. This argument is
pure finite case analysis and imposes no complexity restriction on that
matrix.
\end{proof}
\end{annexitem}

\begin{annexitem}{annex:p05-rep-certified-proof}{thm:p05-rep-certified-proof}{Why enumeration certificates make proof checking primitive recursive}
\begin{proof}
We spell out the link to the governing checker rather than deciding
membership in the c.e. axiom set. The positive-cons coding
\(C(a,b)=1+\operatorname{pair}_{\N}(a,b)\) has components strictly
smaller than its positive code; at zero both projections default to zero.
Iterating the tail at most the initial code reaches the empty list.
Triangular pairing, bounded searches for its inverse, tail iteration,
length, entries, concatenation and filters are primitive recursive, as
explicitly constructed in \cref{thm:arithmetization-syntax}.
Course-of-values tables on child codes similarly compute type, free
variable, numeral and substitution tests; wrong tags, wrong arities,
wrong child types and zero input trees are rejected, not silently
interpreted as formulas.

For a candidate \(N_{15}(R)\), iterate over the finite row list. Each
row has exactly five fields. Check the rule tag against its finite
twenty-rule table, formula types, parameter types, premise arity and
earlier-row references. Bounded sums count parents to ensure a tree with
exactly one parent for each nonroot and none for the last row. Store each
row's computed open-assumption list in the accumulating table. Each
update concatenates child lists and performs the rule's specified
discharge filters and nonvacuity tests. Universal introduction checks
absence from free variables of all remaining assumptions; existential
elimination uses the actual all-occurrence condition in this calculus.
Equality elimination and quantifier eliminations use the already computed
capture-avoiding substitution; alpha conversion is allowed only by its
explicit alpha rule. These are primitive recursive updates iterated for
the primitive recursive row length. Numerical sizes of intermediate
lists need not have a polynomial bound: primitive recursion on the
iteration counter already constructs them.

For \(N_{16}(d,W)\), require a valid derivation with sentence root and
the specified conclusion code, and require the length of \(W\) to equal
the computed root open-list length. At each open occurrence with formula
code \(a\), evaluate the fixed bounded formula
\(\mathsf{Comp}(\overline{e_T},\bar a,\bar s)\) for its supplied
history \(s\). Its numerical truth characteristic is primitive recursive
by structural recursion on this fixed formula, using term evaluation,
Boolean tests, bounded sums and products as above. It is not a uniform
truth predicate for arbitrary arithmetic formulas. The finite loop over
open leaves is therefore primitive recursive. Wrong outer tags,
malformed records, nonsentence roots or leaves, mismatching conclusions,
wrong witness lengths, or a failed history check return zero. A closed
logical derivation requires the empty witness list and makes no membership
call. Together these rules give a finite primitive recursive description
\(d_T\) of the complete Boolean checker, even if membership in \(T\)
is undecidable. The finite-axiom direct-check convention is not substituted
for it halfway through an argument.

The adequacy proof in \cref{thm:arithmetization-syntax} applies unchanged.
For clarity, its two directions use external induction over the finite
postorder rows to identify the checked rule tree and its exact remaining
open occurrences. A successful axiom history means that
\(\varphi_{e_T}(a)\) is defined by \cref{thm:bounded-verification}, hence
that leaf is a given axiom. The extracted tree is a genuine derivation
from \(T\). Conversely take such a finite derivation, put its nodes in
postorder and encode its actual rule parameters and discharge markers.
For each of its finitely many open axiom occurrences, membership in the
domain of \(e_T\) supplies an accepting history by the other direction
of bounded verification. Their finite ordered list supplies \(W\);
none is required for discharged assumptions. All tests then pass.

Apply the characteristic-graph corollary to this particular \(d_T\).
It yields the literal formula \(\exists s\kappa_T(p,x,s)\), its
standard-model accuracy and every positive and negative numeral
instance. Combine graph accuracy with the preceding adequacy directions
to obtain the theoremhood equivalence. The finite aggregate of negative
instances follows from the proved finite-case equivalence for a fixed
\(\bar n\), with this \(\Sigma^0_1\) matrix. Consistency and soundness
have not entered any step. Finally, this selection supplies no
\(\mathsf Q\) equivalence to some other formula solely because the
other formula has the same numeral instances.
\end{proof}
\end{annexitem}

\begin{annexitem}{annex:p05-rep-diagonal-graph}{thm:p05-rep-diagonal-graph}{Total syntax operations, defaults, and strong diagonal representation}
\begin{proof}
The governing constructor codes have proper children smaller than their
parent. Their primitive recursive type and variable tests were proved by
an accumulating course-of-values table in
\cref{thm:arithmetization-syntax}. The numeral code \(\nu(n)\) is
primitive recursion on \(n\). Numeral substitution traverses proper
children, replaces exactly the designated variable occurrences, leaves
other variables and constants unchanged, and commutes with function,
equality and Boolean constructors. At a binder for the designated
variable it leaves the entire subtree untouched; at any other binder it
recurses on the matrix. The inserted numeral has no variables, so no
fresh-variable renaming is needed and no capture is possible. An
accumulating child-code table computes this operation primitively
recursively even though the resulting numeral and formula codes can be
larger than the original code.

For arbitrary term substitution used by the proof checker, the separate
bounded stack traversal in the same coding theorem handles the fresh
index case: take the least index absent from the matrix, replacement
term and designated indices, with bound one above their largest index.
The renaming pass stops at a rebinding of that variable and introduces a
fresh variable absent from the entire matrix, so it needs no further
renaming. The supplied bound \(20(n+1)^2\) on stack updates, where
\(n\) is the original node count, accounts for the main traversal and
its at most \(n\) renaming passes of at most \(n\) nodes. This is a
bounded primitive recursive loop, not an unbounded search. The type test
precedes it and malformed inputs have the declared zero default.

Apply these functions to the input code \(n\), test formulahood and that
the free-variable list is contained in \(\{0\}\), and return its numeral
substitution with \(\nu(n)\) if both tests pass, the fixed code
\(\tau=\ulcorner0=0\urcorner\) otherwise.
Primitive recursive case selection and composition give a fixed
description of this total \(\mathsf{diag}\), including every invalid-code
case. For negation, the same formula type test followed by the single
constructor \(N_7\), or the constant \(N_7(\tau)\) default, gives a primitive recursive
description. Apply the strong-graph theorem separately to each fixed
description. The resulting literal existential-prefix diagonal formula
has a bounded matrix, existence and unrestricted-output uniqueness at
each standard numeral input, exactly as stated. If \(n\) is invalid,
the diagonal biconditional has output \(\bar\tau\), not an unproved partiality exception.
\end{proof}
\end{annexitem}

